%% ma: L11-B16-0054
%% lop: 11
%% chuong: 5
%% bai: 16
%% dang: 11.16.4
%% loai: TN4
%% muc: VD
%% dap_an: "D"
%% nguon: "tai-lieu-hoc-tap-mon-toan-lop-11-theo-chuong-trinh-giao-duc-pho-thong-2018.pdf (Lớp Toán Thầy Duy, Phạm Lê Duy, THPT Chu Văn An), Chương 3 Bài 2, Phần 1, Câu 42 (tr.372)"
%% kien_thuc: "Lớp 11 Bài 16 (giới hạn của hàm số)"
%% trang_thai: chinh-thuc
%% ghi_chu: "Đã duyệt 03/11/2026."
\begin{ex}
Biết rằng $\lim\limits_{x\to-\infty}\left(\sqrt{2x^2-3x+1}+x\sqrt2\right)=\dfrac ab\sqrt2$, ($a;b\in\mathbb Z,\dfrac ab$ tối giản). Tổng $a+b$ có giá trị là
\choice{$1$}{$5$}{$4$}{\True $7$}
\loigiai{\sloppy Nhân liên hợp: $\sqrt{2x^2-3x+1}+x\sqrt2=\dfrac{-3x+1}{\sqrt{2x^2-3x+1}-x\sqrt2}$. Với $x<0$ chia tử và mẫu cho $x$: \[\dfrac{-3+\frac1x}{-\sqrt{2-\frac3x+\frac1{x^2}}-\sqrt2}\to\dfrac{-3}{-2\sqrt2}=\dfrac{3\sqrt2}4.\] Vậy $\dfrac ab=\dfrac34$ và $a+b=7$.}
\end{ex}
